% prezentace.tex -- Lekce 09: Řetězce a radio
% preamble-prez.tex -- sdílená preamble pro prezentace (beamer)
% Includuje se z ucitel/lekceXX/prezentace.tex jako:
%   % preamble-prez.tex -- sdílená preamble pro prezentace (beamer)
% Includuje se z ucitel/lekceXX/prezentace.tex jako:
%   % preamble-prez.tex -- sdílená preamble pro prezentace (beamer)
% Includuje se z ucitel/lekceXX/prezentace.tex jako:
%   \input{../spolecne/preamble-prez.tex}

\documentclass[aspectratio=169]{beamer}

% --- Jazyk a font ---
\usepackage{polyglossia}
\setmainlanguage{czech}
\usepackage{fontspec}

% --- Téma ---
\usetheme{default}
\usecolortheme{default}

\definecolor{microbit-blue}{HTML}{1E4D8C}
\definecolor{microbit-lightblue}{HTML}{D6E4F7}
\definecolor{code-bg}{HTML}{F0F0F0}

\setbeamercolor{frametitle}{bg=microbit-blue, fg=white}
\setbeamercolor{title}{fg=microbit-blue}
\setbeamercolor{block title}{bg=microbit-blue, fg=white}
\setbeamercolor{block body}{bg=microbit-lightblue}
\setbeamercolor{itemize item}{fg=microbit-blue}
\setbeamercolor{itemize subitem}{fg=microbit-blue!70}
\setbeamerfont{frametitle}{size=\large, series=\bfseries}
\setbeamertemplate{navigation symbols}{}
\setbeamertemplate{footline}{%
  \leavevmode%
  \hbox{%
    \begin{beamercolorbox}[wd=.5\paperwidth,ht=2.5ex,dp=1.125ex,leftskip=6pt]{author in head/foot}%
      \usebeamerfont{author in head/foot}\textcolor{gray}{\small Úvod do Pythonu na Micro:bitu}
    \end{beamercolorbox}%
    \begin{beamercolorbox}[wd=.5\paperwidth,ht=2.5ex,dp=1.125ex,rightskip=6pt,right]{date in head/foot}%
      \usebeamerfont{date in head/foot}\textcolor{gray}{\small\insertframenumber\,/\,\inserttotalframenumber}
    \end{beamercolorbox}}%
}

% --- Česky korektní uvozovky ---
\usepackage{csquotes}
\newcommand{\uv}[1]{\enquote{#1}}

% --- Zdrojový kód (Python) ---
\usepackage{listings}
\lstset{
  language=Python,
  basicstyle=\ttfamily\small,
  backgroundcolor=\color{code-bg},
  frame=single,
  framerule=0pt,
  xleftmargin=6pt,
  xrightmargin=6pt,
  breaklines=true,
  showstringspaces=false,
  keywordstyle=\color{microbit-blue}\bfseries,
  stringstyle=\color{teal},
  commentstyle=\color{gray}\itshape,
  literate=
    {á}{{\'a}}1 {é}{{\'e}}1 {í}{{\'i}}1 {ó}{{\'o}}1 {ú}{{\'u}}1
    {ů}{{\r{u}}}1 {č}{{\v{c}}}1 {š}{{\v{s}}}1 {ž}{{\v{z}}}1
    {Á}{{\'A}}1 {É}{{\'E}}1 {Í}{{\'I}}1 {Ó}{{\'O}}1 {Ú}{{\'U}}1
    {Č}{{\v{C}}}1 {Š}{{\v{S}}}1 {Ž}{{\v{Z}}}1 {ň}{{\v{n}}}1 {ř}{{\v{r}}}1,
}

% --- Zvýraznění pojmů ---
\newcommand{\pojem}[1]{\textbf{\textcolor{microbit-blue}{#1}}}
\newcommand{\pyinline}[1]{\texttt{#1}}

% --- Grafika a diagramy ---
\usepackage{tikz}

% --- Ostatní ---
\usepackage{graphicx}
\usepackage{hyperref}
\hypersetup{colorlinks=true, urlcolor=microbit-blue}


\documentclass[aspectratio=169]{beamer}

% --- Jazyk a font ---
\usepackage{polyglossia}
\setmainlanguage{czech}
\usepackage{fontspec}

% --- Téma ---
\usetheme{default}
\usecolortheme{default}

\definecolor{microbit-blue}{HTML}{1E4D8C}
\definecolor{microbit-lightblue}{HTML}{D6E4F7}
\definecolor{code-bg}{HTML}{F0F0F0}

\setbeamercolor{frametitle}{bg=microbit-blue, fg=white}
\setbeamercolor{title}{fg=microbit-blue}
\setbeamercolor{block title}{bg=microbit-blue, fg=white}
\setbeamercolor{block body}{bg=microbit-lightblue}
\setbeamercolor{itemize item}{fg=microbit-blue}
\setbeamercolor{itemize subitem}{fg=microbit-blue!70}
\setbeamerfont{frametitle}{size=\large, series=\bfseries}
\setbeamertemplate{navigation symbols}{}
\setbeamertemplate{footline}{%
  \leavevmode%
  \hbox{%
    \begin{beamercolorbox}[wd=.5\paperwidth,ht=2.5ex,dp=1.125ex,leftskip=6pt]{author in head/foot}%
      \usebeamerfont{author in head/foot}\textcolor{gray}{\small Úvod do Pythonu na Micro:bitu}
    \end{beamercolorbox}%
    \begin{beamercolorbox}[wd=.5\paperwidth,ht=2.5ex,dp=1.125ex,rightskip=6pt,right]{date in head/foot}%
      \usebeamerfont{date in head/foot}\textcolor{gray}{\small\insertframenumber\,/\,\inserttotalframenumber}
    \end{beamercolorbox}}%
}

% --- Česky korektní uvozovky ---
\usepackage{csquotes}
\newcommand{\uv}[1]{\enquote{#1}}

% --- Zdrojový kód (Python) ---
\usepackage{listings}
\lstset{
  language=Python,
  basicstyle=\ttfamily\small,
  backgroundcolor=\color{code-bg},
  frame=single,
  framerule=0pt,
  xleftmargin=6pt,
  xrightmargin=6pt,
  breaklines=true,
  showstringspaces=false,
  keywordstyle=\color{microbit-blue}\bfseries,
  stringstyle=\color{teal},
  commentstyle=\color{gray}\itshape,
  literate=
    {á}{{\'a}}1 {é}{{\'e}}1 {í}{{\'i}}1 {ó}{{\'o}}1 {ú}{{\'u}}1
    {ů}{{\r{u}}}1 {č}{{\v{c}}}1 {š}{{\v{s}}}1 {ž}{{\v{z}}}1
    {Á}{{\'A}}1 {É}{{\'E}}1 {Í}{{\'I}}1 {Ó}{{\'O}}1 {Ú}{{\'U}}1
    {Č}{{\v{C}}}1 {Š}{{\v{S}}}1 {Ž}{{\v{Z}}}1 {ň}{{\v{n}}}1 {ř}{{\v{r}}}1,
}

% --- Zvýraznění pojmů ---
\newcommand{\pojem}[1]{\textbf{\textcolor{microbit-blue}{#1}}}
\newcommand{\pyinline}[1]{\texttt{#1}}

% --- Grafika a diagramy ---
\usepackage{tikz}

% --- Ostatní ---
\usepackage{graphicx}
\usepackage{hyperref}
\hypersetup{colorlinks=true, urlcolor=microbit-blue}


\documentclass[aspectratio=169]{beamer}

% --- Jazyk a font ---
\usepackage{polyglossia}
\setmainlanguage{czech}
\usepackage{fontspec}

% --- Téma ---
\usetheme{default}
\usecolortheme{default}

\definecolor{microbit-blue}{HTML}{1E4D8C}
\definecolor{microbit-lightblue}{HTML}{D6E4F7}
\definecolor{code-bg}{HTML}{F0F0F0}

\setbeamercolor{frametitle}{bg=microbit-blue, fg=white}
\setbeamercolor{title}{fg=microbit-blue}
\setbeamercolor{block title}{bg=microbit-blue, fg=white}
\setbeamercolor{block body}{bg=microbit-lightblue}
\setbeamercolor{itemize item}{fg=microbit-blue}
\setbeamercolor{itemize subitem}{fg=microbit-blue!70}
\setbeamerfont{frametitle}{size=\large, series=\bfseries}
\setbeamertemplate{navigation symbols}{}
\setbeamertemplate{footline}{%
  \leavevmode%
  \hbox{%
    \begin{beamercolorbox}[wd=.5\paperwidth,ht=2.5ex,dp=1.125ex,leftskip=6pt]{author in head/foot}%
      \usebeamerfont{author in head/foot}\textcolor{gray}{\small Úvod do Pythonu na Micro:bitu}
    \end{beamercolorbox}%
    \begin{beamercolorbox}[wd=.5\paperwidth,ht=2.5ex,dp=1.125ex,rightskip=6pt,right]{date in head/foot}%
      \usebeamerfont{date in head/foot}\textcolor{gray}{\small\insertframenumber\,/\,\inserttotalframenumber}
    \end{beamercolorbox}}%
}

% --- Česky korektní uvozovky ---
\usepackage{csquotes}
\newcommand{\uv}[1]{\enquote{#1}}

% --- Zdrojový kód (Python) ---
\usepackage{listings}
\lstset{
  language=Python,
  basicstyle=\ttfamily\small,
  backgroundcolor=\color{code-bg},
  frame=single,
  framerule=0pt,
  xleftmargin=6pt,
  xrightmargin=6pt,
  breaklines=true,
  showstringspaces=false,
  keywordstyle=\color{microbit-blue}\bfseries,
  stringstyle=\color{teal},
  commentstyle=\color{gray}\itshape,
  literate=
    {á}{{\'a}}1 {é}{{\'e}}1 {í}{{\'i}}1 {ó}{{\'o}}1 {ú}{{\'u}}1
    {ů}{{\r{u}}}1 {č}{{\v{c}}}1 {š}{{\v{s}}}1 {ž}{{\v{z}}}1
    {Á}{{\'A}}1 {É}{{\'E}}1 {Í}{{\'I}}1 {Ó}{{\'O}}1 {Ú}{{\'U}}1
    {Č}{{\v{C}}}1 {Š}{{\v{S}}}1 {Ž}{{\v{Z}}}1 {ň}{{\v{n}}}1 {ř}{{\v{r}}}1,
}

% --- Zvýraznění pojmů ---
\newcommand{\pojem}[1]{\textbf{\textcolor{microbit-blue}{#1}}}
\newcommand{\pyinline}[1]{\texttt{#1}}

% --- Grafika a diagramy ---
\usepackage{tikz}

% --- Ostatní ---
\usepackage{graphicx}
\usepackage{hyperref}
\hypersetup{colorlinks=true, urlcolor=microbit-blue}


\title{Lekce 09: Řetězce a radio}
\subtitle{Úvod do Pythonu na Micro:bitu}
\date{}

\begin{document}

\begin{frame}
  \titlepage
\end{frame}

% ------------------------------------------------------------------
\begin{frame}[fragile]{Řetězce — co už umíme}

  Řetězec je text uzavřený do uvozovek:

  \bigskip
  \begin{lstlisting}
zprava = "Hello, Microbit!"

len(zprava)         # 16  -- pocet znaku
zprava + " :-)"     # spojeni (konkatenace)
str(42) + " stupnu" # cislo -> retezec -> spojeni
  \end{lstlisting}

  \pause\bigskip
  Nové operace:

  \begin{lstlisting}
zprava.upper()  # "HELLO, MICROBIT!"
zprava.lower()  # "hello, microbit!"
zprava[0]       # "H"   -- prvni znak
zprava[-1]      # "!"   -- posledni znak
"bit" in zprava # True  -- obsahuje podretezec?
  \end{lstlisting}

  \bigskip
  Indexování funguje \textbf{stejně jako u seznamů} — řetězec je sekvence znaků.
\end{frame}

% ------------------------------------------------------------------
\begin{frame}[fragile]{Řetězcové metody — přehled}

  \begin{center}
  \begin{tabular}{ll}
    \toprule
    \textbf{Výraz} & \textbf{Výsledek (pro \texttt{"Hello"})} \\
    \midrule
    \pyinline{len("Hello")}         & \pyinline{5} \\
    \pyinline{"Hello".upper()}      & \pyinline{"HELLO"} \\
    \pyinline{"Hello".lower()}      & \pyinline{"hello"} \\
    \pyinline{"Hello"[0]}           & \pyinline{"H"} \\
    \pyinline{"Hello"[-1]}          & \pyinline{"o"} \\
    \pyinline{"ell" in "Hello"}     & \pyinline{True} \\
    \pyinline{"xyz" in "Hello"}     & \pyinline{False} \\
    \bottomrule
  \end{tabular}
  \end{center}

  \pause\bigskip
  \begin{block}{Metody vs. funkce}
    \pyinline{len(s)} je funkce — volíme ji s hodnotou jako argumentem.\\
    \pyinline{s.upper()} je \pojem{metoda} — voláme ji \textbf{na řetězci} pomocí tečky.\\
    Obojí jsme ale viděli i dříve: \pyinline{display.scroll()} je také metoda.
  \end{block}
\end{frame}

% ------------------------------------------------------------------
\begin{frame}[fragile]{Radio — bezdrátová komunikace}

  Micro:bit má vestavěný radio čip — posíláme a přijímáme zprávy:

  \bigskip
  \begin{lstlisting}
from microbit import *
import radio

radio.on()              # zapne radio (nutne!)
radio.config(channel=7) # vyber kanal 0-83
  \end{lstlisting}

  \bigskip
  \begin{columns}[T]
    \begin{column}{0.46\textwidth}
      \begin{block}{Odesílání}
        \begin{lstlisting}
radio.send("Ahoj!")
        \end{lstlisting}
        Odešle řetězec všem na stejném kanálu.
      \end{block}
    \end{column}
    \begin{column}{0.46\textwidth}
      \begin{block}{Příjem}
        \begin{lstlisting}
msg = radio.receive()
if msg is not None:
    display.scroll(msg)
        \end{lstlisting}
        Vrátí zprávu nebo \pyinline{None}.
      \end{block}
    \end{column}
  \end{columns}

  \pause\bigskip
  \begin{alertblock}{Stejný kanál}
    Oba micro:bity musí mít \textbf{stejný kanál} — jinak se neslyší.
  \end{alertblock}
\end{frame}

% ------------------------------------------------------------------
\begin{frame}[fragile]{Radio ping-pong}

  Jednoduchý program pro komunikaci dvou micro:bitů:

  \bigskip
  \begin{lstlisting}
from microbit import *
import radio

radio.on()
radio.config(channel=7)

while True:
    if button_a.was_pressed():
        radio.send("Ping!")
        display.show(Image.ARROW_E)

    msg = radio.receive()
    if msg is not None:
        display.scroll(msg)

    sleep(100)
  \end{lstlisting}

  \bigskip
  Oba micro:bity mají \textbf{stejný program}.
  A vidí zprávy \textbf{od všech} na tomtéž kanálu — radio je broadcast!
\end{frame}

% ------------------------------------------------------------------
\begin{frame}{Co jsme se naučili}
  \begin{itemize}
    \item[\checkmark] \pyinline{len(s)}, \pyinline{s.upper()}, \pyinline{s.lower()} — základní operace s řetězci
    \item[\checkmark] \pyinline{s[i]} — přístup k~jednotlivým znakům (indexování jako u seznamů)
    \item[\checkmark] \pyinline{"text" in s} — test, zda řetězec obsahuje podřetězec
    \item[\checkmark] \pyinline{import radio} + \pyinline{radio.on()} — zapnutí rádia
    \item[\checkmark] \pyinline{radio.config(channel=7)} — nastavení kanálu
    \item[\checkmark] \pyinline{radio.send("text")} — odeslání zprávy
    \item[\checkmark] \pyinline{radio.receive()} — příjem zprávy (nebo \pyinline{None})
  \end{itemize}
\end{frame}

\end{document}
